\chapter{Свързани разработки}
\label{ch:svarzani}

Тази глава разглежда съществуващите подходи към моделирането на изразителността —
и в частност на динамиката — в музиката от ударни инструменти и позиционира приноса на
настоящата работа спрямо тях.

\section{Системи основани на правила и статистически подходи}
\label{sec:svarzani-pravila}

Най-ранните системи за изразително изпълнение са \emph{базирани на правила}:
детерминистични евристики, извлечени от музикознанието, които изменят времето и
силата на нотите. Класически пример е системата от правила на KTH за музикално
изпълнение~\cite{friberg2006overview}, която прилага набор от правила (напр.\
акцентиране по силните времена) върху партитурата. Такива системи са
интерпретируеми, но статични: те не се учат от данни и трудно улавят
стилистичното разнообразие между жанрове и изпълнители.

В практиката, някои DAW софтуери имплементират
придаването на „човешко'' звучене обикновено чрез добавяне на
случаен шум (нормално-разпределен) към времето и силата на нотата, или до копиране на готови
шаблони от референтно изпълнение — подходи, които не моделират зависимостта на динамиката от
структурата на конкретната партия.

\section{Невронни модели за груув}
\label{sec:svarzani-nevronni}

Съвременната линия изследвания третира груува (ритъма) като задача за
последователно генериране. \eng{GrooVAE} и свързаният подход
\emph{Learning to Groove}~\cite{gillick2019groove} въвеждат вариационен
автоенкодер над барабанни партии, който умее да придава човешко звучене и да
„квантува наобратно'' (\eng{groove transfer}). Тези модели, както и по-общата фамилия
\eng{MusicVAE}~\cite{roberts2018musicvae}, работят върху \emph{фиксирана мрежа
от шестнайсетини ноти}: партията се представя като матрица от „времеви стъпки
$\times$ инструменти'', придружена от канал с микровремеви отмествания (в милисекунден порядък).

Това представяне е удобно за \emph{генериране на време}, но има две слабости за
нашата цел. Първо, шестнайсетичната мрежа е \emph{дуплетна} и изкривява триолни
и неравноделни усещания (суинг, шафъл, латино), при които две триолни ноти могат
да попаднат в една клетка и да бъдат слети или изпуснати — точно в жанровете,
където динамиката е най-изразителна. Второ, тези модели генерират и времето, и
динамиката \emph{съвместно}, докато при разглежданата тук задача за
\emph{предсказване на динамиката} времето е \emph{известен вход}, а не изход.

\section{Позициониране на приноса}
\label{sec:svarzani-prinos}

Настоящата работа се отличава от изброените по няколко начина. (1)~Разглежда
предсказването на \emph{динамиката} като самостоятелна задача, при която времето
е вход, а не съвместен изход. (2)~Именно защото времето не
се генерира, работата изоставя мрежовото представяне и подава на модела
непрекъсната метрична фаза, като така избягва артефактите от квантуване при
триолни и неравноделни усещания. (3)~Налага строго ограничение срещу изтичане на
информация — никое \eng{MIDI velocity} не служи за признак — което прави
избрания модел \emph{неавторегресивен}. (4)~Сравнява интерпретируема
таблична база (LightGBM), последователен трансформър и вероятностни
изходни глави, като оценява не само точността, но и запазването на динамичния
обхват — свойството, което всъщност определя качеството на възстановената
динамика. (5) Създава приставка за популярен DAW софтуер, която позволява
на музиканти да използват обучените модели на практика.
